% Cote 119: offprints, reports, courses, CV and letters (typed and printed part): archive cote 119, site page 75 (PDF page 76).
% offprint. Draft compared: none.
% AI-transcribed by Claude Opus 5.5 (claude-opus-5-5, adaptive thinking, effort high), Message Batches API, batch r119_full (msgbatch_01DwZNm5Y7YsKBWyAN4sFs8w), 2026-11-13; not yet checked by eye.
\scanpage{119/75}
\mspage{( 3 )}

2.3. Un cas important est celui où la gerbe FIB$_0(C_0)$ est algébrique,
i. e. est liée localement par un groupe de type fini. On dit également que la
catégorie tannakienne $C_0$ est algébrique. On peut voir que $C_0$ est algébrique
si et seulement si elle possède un $\otimes$-générateur [$(^9)$, 4]. Il en résulte que
toute catégorie tannakienne est réunion de ses sous-catégories tannakiennes
algébriques donc que sa gerbe est pro-algébrique, i. e. limite projective de
gerbes tannakiennes algébriques. Réciproquement, si $C$ est une catégorie
ind-tannakienne sur $k$ et si sa gerbe FIB$(C)$ est pro-algébrique, $C$ est
équivalente à une catégorie ind-tannakienne Ind$(C_0)$. J'ignore si cette
condition est toujours satisfaite, elle l'est en tout cas si le lien de FIB$(C)$
[voir $(^4)$, chap. III] est représentable par un groupe. Le résultat essentiel
pour démontrer ce qui précède est le suivant :

Théorème 2.4. --- Si $C$ est une catégorie ind-tannakienne sur $k$ dont la
gerbe FIB$(C)$ est algébrique, $C$ est une catégorie localement noethérienne [$(^3)$,
II, 4], et la sous-catégorie pleine de ses objets noethériens est tannakienne.
De plus, $C$ possède un foncteur fibre sur une extension finie de $k$.

On se sert dans la preuve de ce théorème du résultat suivant d'Algèbre
homologique non-abélienne.

Théorème 2.5. --- Pour une catégorie fibrée $E$ sur Sch$_{/k}$, il est équivalent
d'être une gerbe tannakienne algébrique pour la topologie fpqc ou pour la
topologie fppf.

2. Exemples.

3.1. Soit $k$ un corps, $\dot{M}_k$ la catégorie des motifs sur $k$ $(^2)$; c'est une
$\otimes$-catégorie ACU pseudo-abélienne $\mathbf{Q}$-linéaire, dont la loi $\otimes$ est déduite
du produit direct des $k$-variétés lisses et projectives. Supposons la validité
des conjectures standard [$(^7)$, $(^8)$] pour une théorie de la cohomologie à
valeurs dans un corps K de caractéristique o (par exemple, la cohomologie
$l$-adique). Il en résulte en particulier que $\dot{M}_k$ est munie d'une $\otimes$-graduation
de type $\mathbf{Z}$ : chaque motif M se décompose de façon canonique en une somme
finie
\[ \mathrm{M} = \bigoplus_{n\in\mathbf{Z}} \mathrm{M}^n \]
et, de plus, on a un isomorphisme canonique $(\mathrm{M}\otimes\mathrm{N})^n \simeq \bigoplus_{p+q=n} \mathrm{M}^p\otimes\mathrm{N}^q$.
Modifions la contrainte de commutativité dans $\dot{M}_k$ : si M, N sont des motifs,
et l'isomorphisme de commutativité $\dot{\psi} : \mathrm{M}\otimes\mathrm{N} \xrightarrow{\sim} \mathrm{N}\otimes\mathrm{M}$ a des composantes
$\dot{\psi}^{p,q} : \mathrm{M}^p\otimes\mathrm{N}^q \xrightarrow{\sim} \mathrm{N}^q\otimes\mathrm{M}^p$, le nouvel isomorphisme de commutativité $\psi$
a des composantes $\psi^{p,q} = (-l)^{pq}\,\dot{\psi}^{p,q}$. On vérifie que la nouvelle $\otimes$-caté-%
gorie ACU obtenue, notée $M_k$, est une catégorie tannakienne sur $\mathbf{Q}$. Les
foncteurs fibre sur $M_k$ à valeurs dans des extensions K de $\mathbf{Q}$ sont les théories
de la cohomologie à valeurs dans K; elles vérifient automatiquement les
conjectures standard.

\mspage{\hand{75}}
